\section{Mass Spectrometry Source Settings}
The MS analysis is performed on the Vertex Triple Quad 6500 mass spectrometer equipped with a Turbo V Ion Source operating in electrospray ionisation negative (ESI-) mode. Source parameters were optimized by direct infusion of standard analytes at a concentration of $100~\text{ng/mL}$ at a flow rate of $10~\mu$L/min blended with the mobile phase flow.

\noindent
\begin{minipage}[t]{0.48\textwidth}
    \begin{tcolorbox}[msbox, title=Ion Source Gas Parameters]
        \begin{tabularx}{\textwidth}{X r}
            \textbf{Source Parameter} & \textbf{Value} \\
            \midrule
            Ionisation Mode & ESI Negative \\
            Curtain Gas (CUR) & $35.0~\text{psi}$ \\
            Collision Gas (CAD) & Medium \\
            Ion Source Gas 1 (GS1) & $50.0~\text{psi}$ \\
            Ion Source Gas 2 (GS2) & $60.0~\text{psi}$ \\
        \end{tabularx}
    \end{tcolorbox}
\end{minipage}
\hfill
\begin{minipage}[t]{0.48\textwidth}
    \begin{tcolorbox}[msaccentbox, title=Thermal \& Voltage Settings]
        \begin{tabularx}{\textwidth}{X r}
            \textbf{Source Parameter} & \textbf{Value} \\
            \midrule
            Ion Spray Voltage (IS) & $-4500$~V \\
            Source Temp (TEM) & $450^\circ\text{C}$ \\
            Entrance Potential (EP) & $-10.0$~V \\
            Deflector Potential (DP) & Optimized \\
            Collision Cell Exit (CXP) & $-15.0$~V \\
        \end{tabularx}
    \end{tcolorbox}
\end{minipage}
\par
\vspace{0.5cm}
\subsection{Ionisation Optimization Notes}
Special attention must be paid to the source temperature and spray voltage settings. A higher temperature ($450^\circ\text{C}$) is necessary to ensure rapid and complete desolvation of the droplets, reducing chemical noise and source contamination. The entrance potential is fixed at $-10.0$~V for all transitions to maintain stability of the ion beam before entering the first quadrupole (Q1) \cite{massspec2021}.
